\documentclass[11pt]{article}
\usepackage[T1]{fontenc}
\usepackage{lmodern,amsmath,amssymb,amsthm,mathtools}
\usepackage[margin=1.05in]{geometry}
\usepackage[hidelinks]{hyperref}
\newtheorem{theorem}{Theorem}[section]
\newtheorem{lemma}[theorem]{Lemma}
\newtheorem{proposition}[theorem]{Proposition}
\newtheorem{corollary}[theorem]{Corollary}
\newcommand{\R}{\mathbb R}\newcommand{\E}{\mathbb E}
\newcommand{\D}{\mathcal D}\newcommand{\G}{\mathcal G}
\newcommand{\Hh}{\mathcal H}\newcommand{\code}[1]{{\small\nolinkurl{#1}}}
\title{Time integration and adjoint-graph differentiation\\for reflected-clock Brownian observables}
\author{Project Stardust}
\date{Reviewed editorial draft, October 29, 2026}
\begin{document}
\maketitle
\begin{abstract}
For independent Brownian motions $B,W$, let $Z_t=B(|W_t|)$ and
$A_T=\int_0^T Z_s\,ds$. Although $Z_T$ is outside the original clock
Gaussian Sobolev domain at every positive time, $A_T$ belongs to the
intersection of the two original Gaussian derivative domains. A folded
Gaussian transition bound gives a uniform finite clock energy for spatially
mollified averages, and closedness gives the actual average-domain result.
Bounded smooth functions of averages form a dense derivative domain in
each fixed raw observed past. For the resulting closed derivative and its
maximal Hilbert adjoint, the embedded gradient of $A_T$ is differentiable
in adjoint-graph-dual norm, including left differentiation at the past
endpoint. Its derivative pairs with actual $Z_T$ but has no native field
representative. Along every vanishing positive or negative time-step
sequence, the native gradient-slope norms diverge. These results give
precise source-valued operator identities, not an observed-filtration
It\^o or Clark--Ocone formula. They reconstruct frozen formal results;
novelty remains unresolved.
\end{abstract}

\section{The process, operator and statement}
Let $B$ be a two-sided standard Brownian scenery and $W$ an independent
one-sided standard Brownian motion, on their product Gaussian probability
space $(\Omega,\mathcal F,P)$. Put
\[
 Z_t=B(|W_t|),\qquad A_T=\int_0^T Z_s\,ds\quad(t,T\geq0).
\]
Only the positive scenery branch is used by $Z$. Extending it independently
to negative arguments permits spatial convolution. This is the simple
Brownian-time Brownian model of~\cite{AZ}; an excursion does not select a
fresh scenery, and the clock is reflected. The signed-clock process
$B(W_t)$ is a different model.

Let $D_B,D_W$ be the original closed Gaussian gradients on the product
space, acting respectively in the scenery and clock factors. Their
coefficient Hilbert spaces can be realized as $H_B=L^2(\R)$ and
$H_W=L^2(\R_+)$, using white noise. Equivalently one can use the frozen
library's isometric separable Gaussian coordinates. Additional unused
independent coordinates tensor the graph with constants; averaging over
them gives the reverse domain implication. All fields below retain the
ambient Gaussian codomain
\[
 \G=L^2(P;H_B)\oplus L^2(P;H_W).
\]
They are not assumed observable or predictable in physical time.

Set $E=L^2(\sigma(Z_s:s\geq0),P)$, understood as the closed subspace of
$L^2(P)$ with measurable representatives, and define the \emph{native}
operator
\begin{equation}\label{eq:domain}
 \D f=(D_Bf,D_Wf),\qquad
 D(\D)=E\cap D(D_B)\cap D(D_W).
\end{equation}
The domain is the entire stated intersection. For fixed $H\geq0$ let
\[
 \mathcal F_H^Z=\sigma(Z_s:0\leq s\leq H),\qquad
 E_H=L^2(\mathcal F_H^Z,P),\qquad \D_H=\D|_{E_H}.
\]
This is the raw observed past; $L^2$ always identifies almost-everywhere
equal variables. Taking that quotient does not assert right continuity
of the filtration.

For a closed densely defined operator $T:E_0\supset D(T)\to\G$, write
\[
 \mathcal T_T=D(T^*),\quad
 \|v\|_{\mathcal T_T}^2=\|v\|_{\G}^2+\|T^*v\|_{E_0}^2,
\]
and define continuous linear functionals
\begin{equation}\label{eq:dual}
 (j_TV)(v)=\langle V,v\rangle,\qquad
 (d_Tf)(v)=\langle f,T^*v\rangle.
\end{equation}
The adjoint domain here is maximal: $v\in D(T^*)$ exactly when
$f\mapsto\langle v,Tf\rangle$ is bounded in the $E_0$ norm on $D(T)$.

\begin{theorem}\label{thm:main}
The operators $\D$ and $\D_H$ are closed and densely defined. For every
$T\geq0$, $A_T\in D(\D)$; if $T\leq H$, also $A_T\in D(\D_H)$ and
$\D_H A_T=\D A_T$. Write $V_T=\D A_T$. For every $T>0$,
\begin{equation}\label{eq:fulltime}
 \frac{d}{dT}j_{\D}V_T=d_{\D}Z_T
 \quad\hbox{in }\mathcal T_{\D}'\hbox{ norm}.
\end{equation}
For $0<T\leq H$ and every $h_n<0$ with $h_n\to0$,
\begin{equation}\label{eq:left}
 j_{\D_H}\frac{V_{T+h_n}-V_T}{h_n}
 \longrightarrow d_{\D_H}Z_T
 \quad\hbox{in }\mathcal T_{\D_H}'\hbox{ norm},
\end{equation}
where the finitely many terms with $T+h_n<0$ may be omitted. The limit is
nonzero and has no representative $j_{\D_H}V$ with $V\in\G$. Moreover
\begin{equation}\label{eq:escape}
 \left\|\frac{V_{T+h_n}-V_T}{h_n}\right\|_{\G}\longrightarrow\infty.
\end{equation}
The full operator has the same nonrepresentability, and
\eqref{eq:escape} also holds for every positive vanishing step sequence.
For each $v\in D(\D_H^*)$ the fixed-past derivative has the actual pairing
\begin{equation}\label{eq:actual}
 (d_{\D_H}Z_T)(v)=\E[Z_T\,\D_H^*v].
\end{equation}
\end{theorem}

The proof separates average-domain membership, density on the observed
past, and the change of topology needed for time differentiation. None
of these steps assumes a Gaussian law for $Z$ or a predictable
representation property for its observed filtration.

\section{Time continuity and scenery differentiation}
Let $e_r=\mathbf1_{(0,r]}\in H_B$ for $r\geq0$. The scenery white-noise
isometry sends $e_r$ to $B(r)$, and
\[
 \|B(r)-B(q)\|_{L^2(B)}^2=|r-q|\quad(r,q\geq0).
\]
Conditioning on the entire clock gives
\begin{equation}\label{eq:continuity}
 \|Z_t-Z_s\|_2^2
 =\E\bigl||W_t|-|W_s|\bigr|
 \leq\E|W_t-W_s|
 =\sqrt{2/\pi}\,|t-s|^{1/2}.
\end{equation}
Consequently $Z$ is continuous as an $L^2(P)$-valued curve. It is Bochner
integrable on compact time intervals, and the integral agrees almost
surely with the integral of its continuous paths, by Fubini. The
fundamental theorem of calculus gives
\begin{equation}\label{eq:FTC}
 A'_T=Z_T\quad\hbox{in }L^2(P),\qquad T>0.
\end{equation}
The same derivative holds in the closed subspaces $E$ and, from the
left for $T\leq H$, $E_H$.

For fixed clock sample, $B(|W_t|)$ is first Gaussian chaos in $B$, with
scenery gradient $e_{|W_t|}$. This gives the product closed graph and
\[
 \|D_BZ_t\|_{L^2(P;H_B)}^2=\E|W_t|=\sqrt{2t/\pi}.
\]
Here graph membership is genuine: approximate $e_r$ by finite orthogonal
projections in the scenery coefficient space, and use dominated
convergence with majorant $|W_t|$ for the squared errors. The finite
linear Gaussian graphs converge strongly to the stated product graph.
Integrating that graph yields
\begin{equation}\label{eq:sourceaverage}
 A_T\in D(D_B),\qquad
 D_BA_T=\int_0^T e_{|W_s|}\,ds.
\end{equation}
Indeed both value and gradient are Bochner integrable: the gradient norm
is a constant times $s^{1/4}$, which is integrable on $[0,T]$.
A closed linear graph contains the integral of any integrable curve
in that graph, by approximation with simple graph-valued functions.

\section{The folded transition estimate and average clock domain}
Choose the frozen nonnegative smooth probability mollifiers $\rho_n$,
supported in $[-\epsilon_n,\epsilon_n]$, $\epsilon_n=2/(n+1)$, and let
\[
 B_n(r)=\int\rho_n(u)B(r-u)\,du,\quad
 Z_t^{(n)}=B_n(|W_t|),\quad A_T^{(n)}=\int_0^T Z_s^{(n)}\,ds.
\]
The Brownian increment metric and Minkowski's inequality give, uniformly
in $r\in\R$,
\[
 \|B_n(r)-B(r)\|_{L^2(B)}
 \leq\int\rho_n(u)\sqrt{|u|}\,du\leq\sqrt{\epsilon_n}.
\]
Independence therefore gives
\begin{equation}\label{eq:averageerror}
 \|A_T^{(n)}-A_T\|_2\leq T\sqrt{\epsilon_n}.
\end{equation}
These mollified approximants need not be observed functions of $Z$;
they are used only in the original ambient clock graph.

Put $q_s=\mathbf1_{(0,s]}\in H_W$ and $\operatorname{sgn}(0)=0$.
In the original product Gaussian clock domain,
\begin{equation}\label{eq:clockgraph}
 D_W Z_s^{(n)}=B_n'(|W_s|)\operatorname{sgn}(W_s)q_s,
 \qquad
 G_T^{(n)}=D_W A_T^{(n)}=\int_0^T D_WZ_s^{(n)}\,ds.
\end{equation}
To justify the point graph without assuming the scenery path has a
classical derivative before mollification, expand $B_n(r)$ in a scenery
orthonormal basis. Its scalar coefficients have derivative in the
original one-dimensional Gaussian graph after $r=|W_s|$. The absolute
value has weak derivative $\operatorname{sgn}$; approximating it by
$\sqrt{x^2+\delta^2}$ and then taking the Gaussian closed graph proves
that assertion. For the scenery sums, the square norms of the derivative
coefficients sum to $\|\rho_n\|_2^2$, uniformly in $r$. Thus the value
and derivative sums converge in the product $L^2$ graph. At $s=0$ the
random value is clock independent and the derivative is zero. On $[0,T]$
the derivative norms are bounded by $\sqrt{s}\|\rho_n\|_2$, so graph
integration gives the second equality in \eqref{eq:clockgraph}.

\begin{lemma}\label{lem:folded}
For every $n$ and $T\geq0$,
\begin{equation}\label{eq:energy}
 \|G_T^{(n)}\|_{L^2(P;H_W)}^2
 \leq\frac8{\sqrt{2\pi}}T^{5/2}.
\end{equation}
\end{lemma}
\begin{proof}
By the scenery white-noise isometry,
\[
 \E_B[B_n'(r)B_n'(q)]
 =k_n(r,q):=\int_{\R}\rho_n(r-x)\rho_n(q-x)\,dx\geq0.
\]
For fixed $r$, $\int_{\R}k_n(r,q)\,dq=1$ by Tonelli and normalization.
For $0\leq s<t$, conditional on $W_s=y$, the density of $|W_t|$ on
$q>0$ is
\[
 p_{t-s}(q-y)+p_{t-s}(-q-y)
 \leq\frac2{\sqrt{2\pi(t-s)}}.
\]
Hence $\E k_n(|W_s|,|W_t|)\leq2/\sqrt{2\pi(t-s)}$.
In the double integral for the squared norm of \eqref{eq:clockgraph},
the coefficient inner product is $\langle q_s,q_t\rangle=\min(s,t)$.
After scenery averaging, the signs have absolute value at most one and
$k_n\geq0$. Thus off the time diagonal the absolute integrand is bounded
by
\[
 \frac{2T}{\sqrt{2\pi}}|t-s|^{-1/2}.
\]
For each fixed mollifier, its actual diagonal covariance is finite and
is not assigned zero. The diagonal is a product-Lebesgue null set, so it
does not affect integration. The bound is integrable, since
\[
 \int_0^T|t-s|^{-1/2}\,dt
 =2\sqrt{s}+2\sqrt{T-s}\leq4\sqrt T.
\]
Integrating over $s\in[0,T]$ proves \eqref{eq:energy}. The same domination
justifies the double integrals and conditioning. The constant is a
convenient uniform estimate, not a sharp energy formula.
\end{proof}

\begin{proposition}\label{prop:average}
The actual average $A_T$ belongs to $D(\D)$ for every $T\geq0$. Its
original clock gradient satisfies
$\|D_WA_T\|_2^2\leq8T^{5/2}/\sqrt{2\pi}$.
\end{proposition}
\begin{proof}
For $T=0$ the result is immediate. Otherwise \eqref{eq:averageerror}
gives strong value convergence, and Lemma~\ref{lem:folded} bounds the
actual clock gradients. A bounded sequence in this Hilbert field space
has a weakly convergent subsequence, say $G_T^{(n_k)}\rightharpoonup G$.
The graph of the closed linear operator $D_W$ is a norm-closed subspace,
and hence is weakly closed: a point outside it can be separated by its
orthogonal projection. Therefore
$(A_T^{(n_k)},G_T^{(n_k)})\rightharpoonup(A_T,G)$ lies in that graph.
It follows that $A_T\in D(D_W)$, and lower semicontinuity of the norm
gives the estimate. Equation~\eqref{eq:sourceaverage} gives scenery-domain
membership. Finally $A_T$ is observable from $Z$ and is an element of $E$,
which proves \eqref{eq:domain} membership.
\end{proof}

Uniform derivative bounds alone have been used for membership, not to
assert strong gradient convergence. The library has a further strong
occupation-kernel convergence theorem, but it is unnecessary for the
time-differentiation result proved here.

\section{Density in the fixed raw observed past}
Closedness of $\D$ follows from closedness of both partial graphs and of
$E$. Its restriction $\D_H$ is closed because $E_H$ is closed. Density
requires proof; it does not follow merely from restricting a dense
ambient domain.

\begin{lemma}\label{lem:chain}
Finite linear combinations of $A_t$, $0\leq t\leq H$, belong to
$D(\D_H)$. Their sine and cosine also belong to that domain.
\end{lemma}
\begin{proof}
Linearity and Proposition~\ref{prop:average} prove the first assertion;
all averages up to $H$ are $\mathcal F_H^Z$-measurable. For either
Gaussian factor, the closed gradient satisfies the bounded $C^1$ chain
rule $D\phi(f)=\phi'(f)Df$ when $\phi,\phi'$ are bounded and $\phi'$
is continuous. To see this from the original closure, take smooth
cylindrical graph approximants $f_n\to f$, $Df_n\to Df$. Lipschitz
continuity gives $\phi(f_n)\to\phi(f)$ in $L^2$. The derivative
products converge: the part containing $Df_n-Df$ is bounded by
$\|\phi'\|_\infty\|Df_n-Df\|_2$, and the part containing
$(\phi'(f_n)-\phi'(f))Df$ tends to zero by convergence in probability
and bounded dominated convergence. A subsequence argument gives the
same limit for the whole sequence. Smooth compositions in finitely many
Gaussian coordinates belong to the original closure: cut off the
coordinates, identify the weak partial derivatives by Gaussian
integration by parts, and expand in tensor Hermite polynomials. The
derivative coefficient identities and Parseval show that polynomial
partial sums converge with all first derivatives in $L^2$. Cutoff errors
vanish by Gaussian integrability and the bounded derivative. Thus the
finite-dimensional chain identities hold in the original closed graph,
and closedness proves the asserted chain rule. Apply it in each factor
to sine and cosine. Their
compositions remain observable in $\mathcal F_H^Z$.
\end{proof}

\begin{proposition}\label{prop:dense}
$D(\D_H)$ is dense in $E_H$, including at $H=0$, and $D(\D)$ is dense
in $E$.
\end{proposition}
\begin{proof}
For $0<t\leq H$ and $0<\epsilon<t$, put
$Q_{t,\epsilon}=(A_t-A_{t-\epsilon})/\epsilon$. By
\eqref{eq:continuity}, $Q_{t,\epsilon}\to Z_t$ in $L^2$ as
$\epsilon\downarrow0$. Thus Lemma~\ref{lem:chain} and Lipschitz
continuity show that the sine and cosine of every finite linear
combination of points $Z_{t_1},\ldots,Z_{t_m}$ lie in the $L^2$ closure
of $D(\D_H)$. Zero-time points cause no difficulty since $Z_0=0$.
Constants are in the domain with zero gradient.

Suppose $g\in E_H$ is orthogonal to $D(\D_H)$. It is integrable by
Cauchy--Schwarz. For each finite observation vector
$X=(Z_{t_1},\ldots,Z_{t_m})$, the finite signed measure
\[
 \lambda(S)=\E[g\mathbf1_{\{X\in S\}}]
\]
has all Fourier transforms zero, by the preceding sine/cosine
approximation. Fourier uniqueness implies $\lambda=0$. One can verify
uniqueness by convolving $\lambda$ with a centered Gaussian density:
Fourier inversion for the Gaussian and Fubini make each such convolution
zero; Gaussian approximate identities then recover the integrals of
compactly supported continuous functions, and hence the finite measure.
Therefore $\E[g\mathbf1_A]=0$ for every finite observation cylinder $A$.

These cylinders form a pi-system generating the raw past. The signed
measure $A\mapsto\E[g\mathbf1_A]$ has total mass zero; the pi-lambda
theorem makes it zero on the whole generated sigma-algebra. Since $g$
has a past-measurable representative, taking its positive and negative
sets gives $\E|g|=0$. Hence $g=0$. The orthogonal complement is zero,
which proves density. When $H=0$, $E_0$ consists of constants and the
conclusion is immediate. The same argument with arbitrary finitely
many nonnegative times proves density in $E$.
\end{proof}

This proof uses the actual finite observation laws; it does not replace
them with Gaussian laws or identify conditional expectation with a
source-chaos projection.

\section{Why positive-time points remain outside the domain}
For clarity we recall the clock obstruction underlying the time result,
with its common-field requirement. In the scenery-valued Hilbert space
$K=L^2(B)$, for $c>0$ consider $F_c(x)=B(c|x|)$.
For $0\leq a<x$,
\begin{equation}\label{eq:metric}
 \|F_c(x)-F_c(a)\|_K^2=c(x-a).
\end{equation}
If $F_c$ had a locally square-integrable Bochner weak derivative $g$,
its continuous representative would satisfy
$F_c(x)-F_c(a)=\int_a^x g(u)\,du$. Cauchy--Schwarz and
\eqref{eq:metric} would give
\[
 c\leq\int_a^x\|g(u)\|_K^2\,du.
\]
The right side tends to zero as $x\downarrow a$, a contradiction.

\begin{lemma}\label{lem:point}
For every $T>0$, $Z_T\notin D(D_W)$ and therefore $Z_T\notin D(\D)$.
If $T\leq H$, also $Z_T\notin D(\D_H)$.
\end{lemma}
\begin{proof}
Set $G=W_T/\sqrt T$, a standard Gaussian coordinate, and
$h=\mathbf1_{(0,T]}/\sqrt T\in H_W$. Decompose the remaining clock
Gaussian coordinates orthogonally to $h$. If a full original clock
gradient $V$ existed, contract it with $h$ and conditionally average over
those remaining clock coordinates. Jensen gives a common field
$g\in L^2(\gamma;K)$, where $\gamma$ is standard Gaussian measure.
The product Gaussian integration-by-parts identity, tested against a
scenery variable $e\in K$ and a compactly supported smooth scalar
$p(G)$, gives
\[
 \int\langle F_{\sqrt T}(x),e\rangle(p'(x)-xp(x))\,d\gamma(x)
 =-\int\langle g(x),e\rangle p(x)\,d\gamma(x).
\]
This identity follows first on the Gaussian core and passes to the
closed graph by Cauchy--Schwarz. The recovered $g$ is the same field
for every $e$, a fact essential to the conclusion.

Let $\gamma(x)$ also denote the strictly positive Gaussian density.
Substitute $p=\psi/\gamma$ on any compact interval. Since
$\gamma'=-x\gamma$, one has $(p'-xp)\gamma=\psi'$. The displayed
identity becomes the distributional weak derivative identity relative
to Lebesgue measure. On that interval $g$ is square integrable because
the Gaussian density is bounded below. Testing a countable dense set
of scenery vectors shows that
$F_{\sqrt T}(x)-\int_a^xg(u)\,du$ is one constant $K$-valued vector
almost everywhere. Indeed every scalar projection has zero weak
derivative; separability and a common full-measure set identify that
vector. Both sides have continuous representatives, so the equality
holds throughout the interval. This contradicts \eqref{eq:metric}.
The two native exclusions follow from their domain inclusions.
\end{proof}

This is the point obstruction from the companion~\cite{clock}. The
scenery gradient in Section 2 remains valid; a clock chain rule is not
applied to a point outside its domain.

\section{Adjoint-graph duals and time-slope escape}
We give the closed-operator facts explicitly to fix the test space.
For $T:E_0\supset D(T)\to\G$ closed and densely defined, the maps in
\eqref{eq:dual} satisfy
\begin{equation}\label{eq:bounds}
 \|d_Tf\|_{\mathcal T_T'}\leq\|f\|_{E_0},\quad
 j_T(Tu)=d_Tu,\quad
 \|j_T(Tu)-d_Tf\|_{\mathcal T_T'}\leq\|u-f\|_{E_0}.
\end{equation}
The first bound is Cauchy--Schwarz and $\|T^*v\|\leq\|v\|_{\mathcal T_T}$;
adjointness gives the equality, and applying the first bound to $u-f$
gives the last estimate.

\begin{lemma}\label{lem:representation}
The map $j_T$ is injective, and $d_Tf=j_TV$ for some $V\in\G$ if and
only if $f\in D(T)$, in which case $V=Tf$. If
$u_n\in D(T)$ converges to $f\notin D(T)$ in $E_0$, then
$\|Tu_n\|_{\G}\to\infty$.
\end{lemma}
\begin{proof}
The orthogonal complement of the graph is
\[
 \Gamma(T)^\perp=\{(-T^*v,v):v\in D(T^*)\}.
\]
This follows directly by testing
$\langle x,u\rangle+\langle y,Tu\rangle=0$ on $u\in D(T)$;
density makes $x=-T^*y$ unique. The relation $d_Tf=j_TV$ says precisely
that $(f,V)$ is orthogonal to this complement. Since the graph is
closed, $(f,V)\in\Gamma(T)$. The converse is adjointness. Taking
$f=0$ proves injectivity of $j_T$.

If $\|Tu_n\|$ did not tend to infinity, a subsequence would have bounded
outputs and hence a further weakly convergent subsequence, say to $V$.
Then $(u_n,Tu_n)$ along that subsequence converges weakly to $(f,V)$.
The closed linear graph is weakly closed, so $f\in D(T)$,
a contradiction. This proves full norm escape, rather than mere
unboundedness of one subsequence.
\end{proof}

\begin{proof}[Proof of Theorem~\ref{thm:main}]
Propositions~\ref{prop:average} and \ref{prop:dense}, together with
closedness already noted, give the operator assertions. Put
$V_t=\D A_t$. Adjointness gives $j_{\D}V_t=d_{\D}A_t$. The bounded
linear map $d_{\D}$ commutes with the derivative \eqref{eq:FTC}, proving
\eqref{eq:fulltime}. For fixed $H$, the same equality holds with
$\D_H$ for $0\leq t\leq H$. The left value quotients converge in
$E_H$ to $Z_T$ even when $T=H$. Applying the bound in
\eqref{eq:bounds} proves \eqref{eq:left}; more precisely, its error is
at most
\[
 \left\|\frac{A_{T+h_n}-A_T}{h_n}-Z_T\right\|_2\longrightarrow0.
\]
By Lemmas~\ref{lem:point} and \ref{lem:representation},
$d_{\D_H}Z_T$ is not represented by any $j_{\D_H}V$. It cannot be
zero, because zero has representative $V=0$. The full-operator
conclusion is identical.

The value quotients
$u_n=(A_{T+h_n}-A_T)/h_n$ lie in $D(\D)$ and converge to
$Z_T\notin D(\D)$. Their gradients are
$(V_{T+h_n}-V_T)/h_n$, so Lemma~\ref{lem:representation} proves
\eqref{eq:escape}. The proof applies to positive steps as well.
In particular $V_t$ has no strong ordinary-time derivative at any
positive $T$: such a derivative would force these slope norms to remain
bounded. Finally \eqref{eq:actual} is the definition of $d_{\D_H}$
as an $L^2$ pairing. Its integrand is integrable by Cauchy--Schwarz.
\end{proof}

No tests or past spaces change with the step index in
\eqref{eq:left}. This is differentiation on each \emph{fixed} raw
observed past. It does not resolve convergence of fields with a changing
past, and it is not a temporal stochastic integral identification.

\section{Formal scope and literature}
\begingroup\raggedright
The frozen authority is commit
\code{bab688bb6140cd6316c80680095a914ea0d346db}.
The accompanying theorem map records the exact average graphs, energy
bounds, finite-average chain tests, raw-past density, closed restrictions,
maximal adjoints and time-limit declarations. Its central endpoint is
\code{IBMPastAdjointTimeCalculus}. The full joint operator is
\code{IBMJointIntrinsicDerivative}; its field channels are scenery and
clock in that order. All scalar pairings above are on the actual joint
law, with the Fubini factor order transported isometrically.

The argument retains the original finite-energy domains. The weaker
adjoint-graph dual is explicitly named and its norm is given; it is not
a new native point gradient. The actual full-past adjoint can have a
larger domain than the ambient adjoint restriction, and no equality of
these domains has been assumed. Full baseline build and transitive axiom
evidence were hash-verified; this editorial work creates no new Lean
proof and does not repeat a resource-intensive root build.

The Gaussian derivative/adjoint construction and closed-operator arguments
are classical; see the Gaussian conventions in~\cite{SV}. Allouba and
Zheng~\cite{AZ} supply the Brownian-time model context. Nourdin and
Zeineddine~\cite{NZ} study a signed Brownian clock in two-sided fractional
scenery and an It\^o-type formula, which is a different model and type of
conclusion. The primary model definitions, Gaussian calculus conventions
and relevant introductions were checked. Searches for average-domain
and observed-past analogues were preliminary, not an exhaustive novelty
audit. No independence, priority, or absence-of-precedent claim is made.
The native-domain regularization, fixed-past maximal-adjoint statement
and sharp separation of convergence topologies are the paper's precise
verified content; full physical-time intrinsic calculus remains open.
\par\endgroup

\begin{thebibliography}{9}
\bibitem{AZ} H. Allouba and W. Zheng,
\emph{Brownian-time processes: the PDE connection and the half-derivative
generator}, Ann. Probab. 29 (2001), 1780--1795.
\href{https://arxiv.org/abs/1005.3801}{arXiv:1005.3801}.
\bibitem{SV} T. Sottinen and L. Viitasaari,
\emph{Stochastic analysis of Gaussian processes via Fredholm representation},
Int. J. Stoch. Anal. (2016), article 8694365.
\href{https://doi.org/10.1155/2016/8694365}{doi:10.1155/2016/8694365}.
\bibitem{NZ} I. Nourdin and R. Zeineddine,
\emph{An It\^o's type formula for the fractional Brownian motion in
Brownian time}, \href{https://arxiv.org/abs/1312.0818}{arXiv:1312.0818}
(2013). Model and introduction comparison only.
\bibitem{clock} Project Stardust,
\emph{Gaussian clock-domain obstructions for reflected-clock Brownian
motion}, reviewed editorial companion (2026),
\href{../ibm-clock-domain/main.pdf}{accompanying manuscript}.
\end{thebibliography}
\end{document}
